\section{Introduction}\label{sec:intro}

Large language model (LLM) agents are usually built, deployed, and evaluated as episodic systems. An agent receives a task, reasons over it, calls tools, returns an answer or performs an action, and then resets. Whatever it learned during the episode, which observations it gathered, which tool outputs it trusted, which sub-goals it abandoned, is discarded at the boundary. Under this contract the agent is stateless between tasks by construction, and most of the safety and correctness reasoning a designer must do collapses to a single episode: if the current prompt is clean and the current tool calls are authorized, the agent is well behaved. The episodic contract is convenient precisely because it makes the past irrelevant.

Always-on agents break this contract. An always-on agent is one whose behavior at a given moment depends on state accumulated before that moment, beyond the current prompt or task instance. A coding assistant that remembers a repository's conventions across weeks, a personal assistant that tracks a user's evolving preferences and open commitments, a customer-support agent that carries a case across many sessions and operators, and a fleet of collaborating agents that share a common workspace are all always-on in this sense, regardless of whether they literally run every second. The defining property is not continuous execution but persistence: the agent retains user preferences, task traces, tool outcomes, procedural lessons, policy constraints, social context, and past mistakes, and any of these can shape, or authorize, its next action. This is a qualitative shift, not a quantitative one. Once retained state can license future behavior, the past is no longer irrelevant, and the single-episode safety argument no longer holds.

\paragraph{State becomes authoritative.} Persistence changes what stored information can do. A stored preference can silently fix the default arguments of a tool call. A cached credential can grant access long after the user intended it to lapse. A procedural skill distilled from one successful trajectory can be replayed as an executable commitment in a situation it no longer fits. A summary written months ago can override fresh evidence because it retrieves more strongly. In each case the agent is not reasoning from the current situation alone; it is being guided by a record whose authority, scope, freshness, and origin were fixed in the past and are rarely re-examined. An always-on agent is therefore more than a memory-augmented episodic agent. It is a \emph{persistent-state system} whose retained state needs rules for what may persist, what may act, and what may be revoked.

We use \emph{persistent state} as a broad term. It includes retrievable memory, the object most prior work studies, but also the durable, control-relevant records that surround memory and cross the episode boundary with the agent: task ledgers of open commitments and obligations, permissions and consent, tool and credential state, provenance and audit trails, shared and social state across users and agents, and trigger state governing when to act proactively. External commitments are not internal memory records; they are part of the agent's persistent accountability surface because they must remain traceable to the internal state that authorized them. A calendar reminder, a partially completed purchase, a revoked API token, a poisoned workflow, and a deletion request that has only partly propagated are all persistent-state concerns, even when none is naturally described as a ``memory.'' This wider unit separates the survey from memory-form taxonomies.

\figevolution
\FloatBarrier

\paragraph{Persistence is useful, and persistence is risky.} The case for persistence is strong. Retaining state lets an agent avoid repeating work, personalize its behavior to a particular user or codebase, and adapt over time without retraining or weight updates, simply by accumulating and reusing experience. The same property that makes persistence useful, however, makes it dangerous in ways an episodic agent never faces. Irrelevant or salient-but-wrong memory can distract retrieval and crowd out current evidence. Stale memory can override fresh observation because nothing marks it as expired. Private memory can leak across users or sessions when scope is not enforced. Poisoned memory, written during one ordinary interaction, can persist and activate much later, after the attacker's session has ended and the triggering context is gone. Lossy consolidation can delete the precise identifier, date, or handle that a later action needs, while leaving a summary that still scores well on aggregate recall. None of these failures is repaired by filtering the current prompt, because the harm already resides in durable state. Persistence converts transient, single-episode problems into long-lived, system-level ones, a failure class episodic evaluation is not built to see.

\paragraph{The unit of analysis: persistent state, not memory.} The literature on LLM agents is large and growing, and much of it touches the questions above. Broad agent surveys map perception, reasoning, planning, tool use, and multi-agent organization \citep{xi2023riseandpotential,wang2023surveylargelanguagemodelagents}. Within that space, memory surveys explain how agents remember: the forms memory takes, the operations of writing, organizing, retrieving, and forgetting, and the benchmarks that probe recall \citep{zhang2024memorymechanism,du2026memoryautonomous,hu2025memoryageaiagents}. A more recent line argues that current memory systems are retrofitted from retrieval and database stacks rather than designed for agents, and calls for agent-native memory and externalized agent infrastructure \citep{zhou2026agentnative,zhou2026externalizationllmagentsunified}. We use much of this literature but change the unit of analysis. The first question is not what form an agent's memory takes, or which retrieval mechanism works best. It is how retained state \emph{authorizes} future behavior, how it propagates across sessions, users, and agents, and how it can be \emph{repaired} when it becomes stale, poisoned, overscoped, or wrong. Memory is one component of that object; governance of the whole remains under-specified.

This reframing follows a lineage that predates LLMs. Cognitive science long distinguished episodic from semantic memory \citep{tulving1972episodic}, and cognitive architectures such as Soar and ACT-R made the further distinction between declarative facts retrieved by activation and decay and procedural rules selected by utility, embedding all of them in an explicit decision cycle \citep{laird2012soar,anderson2004actr}. These architectures already treated memory not as a passive store but as a governed component of an acting system, with separation between memory types and explicit operations over them. Modern LLM agents inherit the vocabulary, with CoALA mapping cognitive-architecture structure onto language agents by organizing memory modules and an action space \citep{sumers2024cognitivearchitectureslanguageagents}. What they do not yet inherit is the runtime enforcement of the boundaries those architectures assumed: in an LLM agent, nothing prevents a procedural lesson from silently overwriting a semantic fact, an expired permission from licensing a tool call, or a deleted record from surviving in a derived tier. The gap between a memory \emph{type} and a governed \emph{state record} is the gap this survey is about.

\paragraph{Contributions.} This survey makes four contributions, which also organize the argument:

\begin{enumerate}
\item \textbf{Definition.} We define always-on agents as persistent-state systems and characterize a persistent-state item along six diagnostic axes that can be probed separately in simple cases but often interact in deployed systems: \emph{authority} (what permits this state to influence an action), \emph{scope} (which user, task, tool, time, or group may use it), \emph{mutability} (whether it can be revised, superseded, decayed, or locked), \emph{provenance} (which source, timestamp, and transformation produced it), \emph{recoverability} (whether derived state and affected decisions can be rolled back), and \emph{actionability} (whether it is evidence, preference, policy, skill, or executable commitment). A memory that scores perfectly on recall can still be unsafe if it has no authority boundary, no provenance, and no path to repair after a bad action.
\item \textbf{Lifecycle.} We recast memory as one stage of a broader governed loop, observe, write, validate, organize, retrieve, act, update, forget, audit, and rollback, and name five invariants the loop should preserve but that an episodic agent satisfies vacuously by resetting: authority monotonicity, scope non-expansion, deletion propagation, provenance preservation, and rollback traceability. Each invariant is the property whose violation defines a persistence-specific failure class.
\item \textbf{Coverage map.} We place memory, RAG, long-context, embodied-agent, and social-agent mechanisms and benchmarks against the always-on requirements they report testing, backed by a $435$-work coded corpus. In our coding, the largest qualitative skew is that work concentrates on accumulating and retrieving state more than on governing it.
\item \textbf{Protocol and pilot.} We propose an Always-On Evaluation Protocol that scores state mutation and recovery rather than answer quality alone, and instantiate it in AOEP-v0: an event-stream and snapshot schema, worked episodes, a validator design, and a pilot that illustrates which governance fields current memory wrappers do not expose.
\end{enumerate}

\figoverview

\paragraph{Reader's guide.} The survey is organized as a layered stack over persistent state, with three cross-cutting lenses, sketched in Figure~\ref{fig:overview}. After this introduction, we fix the definition and the six axes, then use the lifecycle and its five invariants as the main frame for the paper. The stack then moves through where state physically lives, from model parameters and the context window through retrieval stores, knowledge graphs, multimodal stores, and operating-system-style serving runtimes (substrates); how state moves through the forward arc as typed transitions (the state lifecycle); how state changes over a long horizon, where the classic plasticity-stability tension reappears in externalized form (continual adaptation); how state is governed along the return arc of update, forget, audit, and rollback (governance); and how state is shared and propagates across agents (multi-agent and shared state). The three cross-cutting parts read the stack through evaluation and benchmarks, failure modes and security, and application domains, before a research agenda and conclusion. Readers interested mainly in the empirical gap can read this introduction and the governance and evaluation parts; readers interested in mechanisms can enter at the substrate and lifecycle parts. Table~\ref{tab:roadmap} maps each section to the regime it serves, the question it answers, and whether it lies on a suggested minimal reading path for a time-constrained reader. The remainder of this introduction makes the corpus and its construction explicit, then positions the survey against its closest neighbors.

\begin{table}[t]
\centering
\small
\setlength{\tabcolsep}{5pt}\renewcommand{\arraystretch}{1.15}
\begin{tabularx}{\linewidth}{p{0.045\linewidth}p{0.20\linewidth}p{0.115\linewidth}Yc}
\toprule
\S & Section & Regime & Question it answers & Core path \\
\midrule
2 & Defining always-on agents & \substance & What is an always-on agent, and what are the six state axes? & $\bullet$ \\
3 & Persistent-state taxonomy & \substance & What kinds of state does the agent hold? & $\bullet$ \\
4 & The state lifecycle & \motion & By what operations does state move, and what invariants govern it? & $\bullet$ \\
5 & State substrates & \substance & Where does state physically live? & \\
6 & Mechanism families & \motion & How do deployed systems write, adapt, and share state? & \\
7 & Evaluation and benchmarks & \control & What does the field measure, and what does it not? & $\bullet$ \\
8 & Failure modes and governance & \control & How does persistence fail, and what governs it? & $\bullet$ \\
9 & The AOEP protocol & \control & How can governance be scored rather than only described? & $\bullet$ \\
10 & Resources and domains & \control & What does each application domain demand of persistence? & \\
11 & Research frontiers & \control & What are the open problems and how might they be solved? & $\bullet$ \\
12 & Conclusion & n/a & What follows for the field? & $\bullet$ \\
\bottomrule
\end{tabularx}
\caption{Reading roadmap. Each section is placed in one of the three regimes the survey is organized around (\substance: what state is and where it lives; \motion: how it moves; \control: how it is governed, judged, and stress-tested). The final column marks a suggested minimal path for a time-constrained reader who wants the empirical argument and the protocol without the full mechanism and substrate catalogues.}
\label{tab:roadmap}
\end{table}

\subsection{Survey scope and corpus}\label{sec:intro:scope}

\paragraph{What counts as in scope.} We study persistent-state mechanisms, the benchmarks and evaluation methods that probe them, the failure modes persistence induces, the foundational work that anchors our taxonomy, and the boundary cases that sharpen the definition. We include a work if it builds or studies a persistent-state mechanism, contributes a relevant benchmark or evaluation method, documents a persistence-induced failure mode, anchors the state taxonomy or the plasticity-stability framing, or defines a boundary case. We exclude pure model-architecture or pretraining work with no persistent-state operation, agent applications that add no persistent-state operation over an off-the-shelf memory, and superseded versions of works we already include. Long-context models, RAG systems, and context-engineered agents sit deliberately near the boundary and are admitted as constraints rather than as the object itself: long-context benchmarks ask whether a model can use what is already in its window, RAG benchmarks ask whether external evidence is selected and used faithfully, and context-engineering work systematizes how retrieval, memory, tools, and multi-agent signals are assembled for one inference \citep{mei2025surveycontextengineering}. Always-on agents add a different question these lines do not pose: which of the assembled signals should become agent-owned longitudinal state, when should that state become authoritative, and how can it later be revised, scoped, revoked, or removed. Personalization is treated the same way, as a state \emph{type} rather than the definition; we return to it below.

\paragraph{Search protocol.} We assembled the corpus by a recorded protocol over arXiv (cs.AI, cs.CL, cs.LG, cs.CR), Semantic Scholar, OpenReview, and the ACL Anthology, covering 2023 to 2026, with earlier foundational anchors admitted when they ground a taxonomy or a term we use. The pre-2023 anchors are deliberately upstream of LLM agents and concentrate in the cognitive and neuroscience foundations: the episodic-semantic distinction \citep{tulving1972episodic}, multi-component working memory with an episodic buffer \citep{baddeley2000episodic}, complementary learning systems and their update for artificial agents \citep{mcclelland1995complementary,kumaran2016complementary}, the neural evidence for offline replay during sleep \citep{wilson1994reactivation}, quantitative forgetting curves \citep{murre2015replication}, and the Soar and ACT-R architectures \citep{laird2012soar,anderson2004actr}. We seeded the modern corpus from agent-memory surveys and a closest-work set, chased citations two hops in each direction, and then ran targeted term-set sweeps. Most sweeps aimed away from the well-covered center of agent memory and toward the periphery we expected to be thin: authority, rollback, formal verification, durable execution, deletion propagation, cross-agent interoperability, and the database, distributed-systems, reinforcement-learning, and HCI literatures that publish persistent-state work outside agent-memory vocabulary. This matters for interpreting the skew below, because the sweeps were designed to \emph{over}-sample governance, not to under-sample it.

\paragraph{Admission flow.} Table~\ref{tab:corpus-flow} makes the corpus build less of a black box. It records the seed set, every admitted non-duplicate addition, and the round structure used to fold new works into the coded corpus. The flow is admission-based rather than a full PRISMA exclusion diagram: later rounds record rejected and duplicate candidates, but the earliest rejected-candidate identifiers were not retained in a uniform screened/excluded table. We therefore use the flow to support directional coverage claims and stopping conditions, not prevalence estimates over all papers that could have been screened.

\begin{table}[t]
\centering
\small
\setlength{\tabcolsep}{4pt}\renewcommand{\arraystretch}{1.16}
\begin{tabularx}{\linewidth}{p{0.18\linewidth}p{0.31\linewidth}p{0.11\linewidth}Y}
\toprule
Corpus step & Admission rule & Added & What it tested \\
\midrule
Seed and closest-work set & Recent memory and agent surveys plus hand-identified closest neighbors satisfying I1 to I5 & 97 & Establishes the vocabulary and the first boundary around persistent-state work \\
Rounds 1 to 2 & Non-duplicate mainstream agent-memory and benchmark additions & 147 & Expands the center of the memory and evaluation literature \\
Rounds 3 to 10 & Targeted governance, rollback, authority, deletion, durable-execution, HCI, and systems sweeps & 171 & Deliberately over-samples the return arc that the survey expects to be sparse \\
Round 11 & Saturation check over mainstream agent-memory vocabulary & 2 & Tests whether the memory frontier still yields many uncoded works \\
Rounds 12 to 15 & Adjacent-disciplinary and confirmation sweeps & 18 & Probes work hidden under database, distributed-systems, formal, and organizational vocabulary \\
\midrule
\textbf{Total} & \textbf{Seed plus fifteen recorded non-duplicate rounds} & \textbf{435} & \textbf{Coded by category, lifecycle stage, state axis, and subarea} \\
\bottomrule
\end{tabularx}
\caption{Admission flow for the coded corpus. Counts are admitted works, not candidate-screening totals. The flow explains how the corpus reached $N{=}435$ and why the governance skew is read as a large pattern inside our query frame rather than as a field-wide census.}
\label{tab:corpus-flow}
\end{table}

\paragraph{A category $\times$ lifecycle-stage $\times$ state-axis $\times$ subarea coding scheme.} Every included work is coded along four dimensions. First, a \emph{category}: mechanism, benchmark, failure-mode, survey, foundation, or boundary. Second, a multi-label set of \emph{lifecycle stages} it primarily exercises, drawn from the ten stages observe, write, validate, organize, retrieve, act, update, forget, audit, and rollback. Third, a multi-label set of \emph{state axes} it touches, drawn from the six axes authority, scope, mutability, provenance, recoverability, and actionability. Fourth, a primary application or method \emph{subarea}, of which thirty-three recur (the largest being conversational memory, memory mechanism, security and poisoning, evaluation methodology, multi-agent social memory, RAG and long-context, multimodal memory, state governance, agentic-OS runtime, continual learning, and skill learning). The cross-product of these dimensions is what lets the corpus answer questions a flat reading list cannot, for example which lifecycle stages the benchmark literature actually tests, or which state axes the mechanism literature leaves implicit. Appendix~\ref{app:methods} gives the full protocol: exact source list and time window, the query term sets, the inclusion and exclusion criteria, the coding scheme, and a measured inter-coder reliability check (pooled agreement $0.82$ on lifecycle stages and $0.74$ on state axes over a blind $236$-work sample), together with the threats to validity that bound how the headline counts should be read.

\paragraph{The corpus and its skew.} The resulting corpus of $435$ works is summarized in Table~\ref{tab:intro-corpus}, with its primary mapping into the survey's nine parts. Every work has a primary part, no part is empty, and the lightest part still hosts more than twenty works. The notable feature of the corpus is not its size but its shape. In our coding, coverage concentrates on the early, accumulative end of the lifecycle: retrieve ($269$ of $435$) and write ($200$) dominate the stages, and mutability ($160$) and provenance ($153$) lead the axes. Coverage thins sharply on the governance end. Authority is the rarest axis at $72$ of $435$. Among lifecycle stages, audit ($88$), forget ($66$), and especially rollback ($27$) trail far behind retrieval. Only twenty-seven of $435$ works expose any rollback mechanism for state-affected decisions, and none in the corpus reports recovery success or cost after corruption. These counts are scoping estimates, not a census, but they point to a large qualitative asymmetry: the works we coded have many techniques for accumulating and recalling state and fewer for governing it. Figure~\ref{fig:timeline} shows the same skew over time: the early-stage share stays flat and high from 2023 to 2026 while the governance share, though rising, remains a minority and turns upward only from 2025.

\begin{table}[t]
\centering
\small
\begin{tabular}{clr}
\toprule
Part & Title & Works \\
\midrule
P1 & Foundations, Definitions, and State Types & 27 \\
P2 & State Substrates and Representations & 117 \\
P3 & State Lifecycle (Write/Organize/Retrieve/Act) & 33 \\
P4 & Continual Adaptation and Plasticity-Stability & 52 \\
P5 & Governance (Authority/Provenance/Deletion/Audit) & 53 \\
P6 & Multi-Agent and Shared State & 23 \\
P7 & Evaluation and Benchmarks & 25 \\
P8 & Failure Modes and Security & 25 \\
P9 & Application Domains & 80 \\
\midrule
\multicolumn{2}{l}{\textbf{Total}} & \textbf{435} \\
\bottomrule
\end{tabular}
\caption{The $435$-work coded corpus, mapped to the survey's nine parts (primary subarea mapping; multi-label cross-references add more). Each work is also coded by category, lifecycle stage, and state axis. The early-lifecycle, low-governance skew described in the text (retrieve $269$, write $200$ versus rollback $27$; authority the rarest axis at $72$) holds within our coding and should be read as a scoped estimate rather than a census of the field.}
\label{tab:intro-corpus}
\end{table}

\paragraph{Governance-targeted search within the query frame.} One possible objection is that we looked in the wrong places, and that the governance literature is larger than our corpus suggests. Our evidence cannot rule out that possibility, but it does constrain the simplest version of it within our recorded query frame. First, the bulk of our search rounds targeted governance, rollback, formal verification, durable execution, deletion propagation, and adjacent systems literatures; those rounds grew the corpus more than fourfold yet only lifted the governance fraction from roughly one-sixth to about one-third before late rounds showed diminishing returns under this query frame, with rollback still rare. Second, the skew worsens, not improves, when we restrict to work submitted through 2024: the governance fraction falls rather than rises, so the corpus pattern is not created solely by including immature 2026 preprints. Third, the skew appears across categories and subareas rather than in a single corner. We therefore read the corpus shape as evidence of a qualitative asymmetry in the works we coded, with the standard caveat that this is representative coding rather than an exhaustive census.

\paragraph{Reliability and diminishing returns.} Two checks bound the confidence we place in these counts. For \emph{reliability}, labels come from a single coding pass with a blind second-coder check on a sample; the multi-label stage and axis dimensions are the harder cases, since a system that writes, retrieves, and updates state legitimately receives all three stage labels, and we treat such overlap as signal rather than noise. For \emph{search coverage}, the late, governance-targeted rounds added many works but did not move the governance fraction or introduce new lifecycle stages or state axes, which suggests diminishing returns within the directions we sampled rather than a proof of exhaustive coverage. Personalization is the one subarea where the recent literature is still expanding rapidly, and we read its newest entries as evidence of emerging directions rather than settled consensus. We treat all figures here as scoping estimates whose directional conclusion is that the works we coded accumulate state more often than they govern, relinquish, or repair it.

\subsection{Related surveys and how this survey differs}\label{sec:intro:related}

The obvious risk is that this reads as a relabeled memory survey. The adjacent survey literature is already dense, so the paper must make the distinction concrete. This survey integrates lifecycle, security, policy, and state-trajectory threads under \emph{persistent state} as the unit of analysis, then attaches a checkable scoring contract to the obligations that follow from that unit.

\paragraph{Agent surveys.} The broadest neighbors survey LLM agents as a whole, mapping perception, reasoning, planning, tool use, and multi-agent organization \citep{xi2023riseandpotential,wang2023surveylargelanguagemodelagents}. These works are indispensable for situating agents. Their native goal is to explain agent architectures and capabilities, so memory is usually one functional module among many rather than the unit whose authority, provenance, and recoverability are tracked. Our survey zooms into the state that such a module, and the records around it, carry across episodes, and asks how that state is governed rather than how the surrounding architecture is organized.

\paragraph{Memory surveys.} The closest and most numerous neighbors survey memory mechanisms for LLM agents directly. They systematize memory \emph{forms} (parametric versus non-parametric, episodic versus semantic versus procedural, short- versus long-term) and memory \emph{operations} (encode, store, retrieve, forget), and they catalog the benchmarks that test recall \citep{zhang2024memorymechanism,du2026memoryautonomous,hu2025memoryageaiagents,huang2026rethinkingmemorymechanisms}. Their native contribution is to explain how agents remember and manage memory. Several already model a write, manage, and read loop \citep{du2026memoryautonomous}, give a broad descriptive taxonomy of forms, dynamics, and benchmarks \citep{hu2025memoryageaiagents}, or ask when memory is useful rather than simply present \citep{huang2026rethinkingmemorymechanisms}. Our survey builds on that base but changes the unit: permissions and credential state, task ledgers and commitments, audit trails, rollback handles, and externally committed effects linked to internal records are not memory forms in the usual sense, but they determine what retained state may govern later action. The delta is therefore an operation-level handle on the state record itself: modeling lifecycle operations as typed transitions over records with authority, scope, provenance, and recoverability links, then asking which obligations those transitions should preserve.

\paragraph{Externalization and agent-native position papers.} A third family argues, correctly in our view, that today's memory systems are often retrofitted from retrieval and database stacks and that agents need native abstractions and externalized infrastructure \citep{zhou2026agentnative,zhou2026externalizationllmagentsunified}. We share the diagnosis and treat it as motivation. Their native goal is architectural: to name the components an agent infrastructure should expose. Our contribution is complementary. We ask how revocation, deletion completeness, rollback, and authority currency become properties of state transitions rather than only desiderata attached to components, and then turn those transition obligations into an evaluation contract.

\paragraph{Long-term-memory security surveys.} A fourth family surveys the security of long-term agent memory, organizing threats and defenses across write, retrieve, forget, audit, and rollback phases \citep{lin2026surveylongtermmemorysecurity}. This is the neighbor closest to our governance and failure-mode parts, and it correctly centers the lifecycle phases where memory is attacked. Its native goal is security taxonomy and defense mapping. Our survey is broader in mechanism, evaluation, substrate, and domain coverage, and narrower in one methodological sense: we ask which persistent-state transitions would have to be logged or checked for a defense claim to become an executable obligation.

\paragraph{Personalization as an adjacent boundary, not the definition.} Personalization deserves separate treatment because it is the adjacent area most often confused with always-on agency, and because its recent literature is the most active. Benchmarks here evaluate user-profile-conditioned outputs and, increasingly, the parts closest to persistent state: following stated preferences across long multi-session context \citep{zhao2025prefeval}, tracking an evolving persona over interactions \citep{jiang2025personamem}, surveys that organize the foundations and evaluation of personalized LLM agents \citep{xu2026personalizedllmpoweredagentsfoundations,salemi-etal-2024-lamp}, and a fast-growing 2025 to 2026 wave that pushes on individual state properties: principled forgetting balanced against recall \citep{uddin2026recallforgettingbenchmarkinglongterm}, write-consolidate-retrieve coordination beyond plain retrieval \citep{li2026himes}, attributes that drift on different timelines \citep{xie2026dynamicmem}, recall of memories logically critical but semantically distant from the query \citep{ding2026imlogic}, reuse of persisted preferences to fill under-specified tool arguments across sessions \citep{yoon2026mpt}, category-bounded long-term memory for an in-car assistant \citep{kirmayr2025carmem}, retrieval of historical web behaviors to align a personalized web agent \citep{cai2025personalizedwebagents}, production-scale industrial memory over noisy longitudinal data \citep{xu2026linkedinmemory}, accumulation of personalized context to resolve implicitly-specified targets \citep{lee2026polar}, long-term personalized assistants that capture subjective user traits \citep{huang2025mempal}, and auditing long-term memory as a post-interaction record rather than only through downstream task success \citep{ma2026memprobe}. This wave is evidence that the field is converging on persistent state from the user-modeling side. Yet personalization remains a different object than ours. A system can top a preference-following benchmark \citep{zhao2025prefeval} while having no mechanism to revoke a preference the user later retracts, and high recall of a preference whose authority has lapsed is not a success but an authority-monotonicity failure these benchmarks do not score. One line makes the danger explicit by formalizing unintended long-term state poisoning, where routine interactions gradually weaken confirmation boundaries and silently expand the agent's action scope \citep{xu2026toxicchats}: an accumulation failure that only a governance frame can name. We therefore treat personalization, like trigger state and the other categories, as a persistent-state \emph{type} that our axes and lifecycle must cover, not as the definition of always-on agency.

\begin{table}[t]
\centering
\scriptsize
\setlength{\tabcolsep}{3pt}\renewcommand{\arraystretch}{1.18}
\begin{tabularx}{\linewidth}{p{0.20\linewidth}p{0.18\linewidth}p{0.19\linewidth}p{0.18\linewidth}Y}
\toprule
Closest work & Unit of analysis & Lifecycle or evaluation coverage & Governance coverage & What this survey adds \\
\midrule
Agent-memory surveys \citep{du2026memoryautonomous,huang2026rethinkingmemorymechanisms,hu2025memoryageaiagents} & Memory mechanisms and utility & Write, manage, retrieve; substrates and cognitive forms & Partial treatment of control, forgetting, and policy & Persistent state as the unit, including permissions, ledgers, credentials, side effects, and recovery obligations \\
Long-term-memory security \citep{lin2026surveylongtermmemorysecurity} & Attacks and defenses over memory phases & Security lifecycle around write, retrieve, forget, audit & Strong threat taxonomy; less complete substrate, evaluation, and domain map & Integrates security with mechanisms, benchmarks, state types, and a transition-level evaluation contract \\
Governed evolving memory \citep{orogat2026gem} & State-trajectory correctness & Formal operators and correctness conditions & Explicit rollback-traceability view & Positions AOEP as a prototype contract that instantiates related obligations in executable form \\
Agent-evaluation surveys \citep{yehudai2025surveyevaluationllmagents,mohammadi2025evalbenchmarkingllmagents} & Agent task and process evaluation & Planning, tool use, memory, self-reflection, enterprise concerns & Authority, deletion, and rollback mostly not scored as state obligations & Asks which persistent-state properties existing agent evaluations operationalize \\
Evaluation methodology \citep{kapoor2024aiagentsthatmatter} & Benchmark quality and reporting & Cost, overfitting, holdout, failure-rate critique & Methodological rather than state-specific & Uses these standards to constrain how AOEP and corpus claims should be reported \\
Tool-agent security \citep{debenedetti2024agentdojodynamicenvironmentevaluate} & Prompt injection in dynamic tool environments & Realistic tool/data attacks and defenses & Strong on untrusted data and tool boundaries; lighter on persistence & Identifies what changes when injected data becomes durable state across sessions \\
Contextual privacy \citep{mireshghallah2025cimemories,wen2026contextualizedprivacy} & Appropriate information flow from memory & Contextual integrity and context-aware privacy defense & Strong scope lens; limited rollback and lifecycle recovery & Folds contextual integrity into scope non-expansion across persistent state transitions \\
\bottomrule
\end{tabularx}
\caption{Closest-neighbor comparison. The novelty claim is integration and operationalization: this survey unifies related threads under persistent state, makes authority and recoverability first-class diagnostic axes, and asks which state-transition obligations can be checked.}
\label{tab:closest-surveys}
\end{table}

\paragraph{The common gap.} Across these neighboring families, the opportunity is integration rather than replacement. Memory surveys taxonomize forms and recall; agent surveys situate memory as a module; externalization papers inventory components; security surveys catalog attacks and defenses; privacy and evaluation work cover important slices of scope, safety, and measurement. What remains under-specified is the governed state record as a common unit: lifecycle operations as transitions over that record, and obligations that can be checked under restart, conflict, deletion, adversarial write, shared-scope change, and delayed consequence. This survey is therefore not another taxonomy of how agents \emph{remember}. It is an account of how agents \emph{persist}, and of why the works we coded accumulate and retrieve state more often than they govern, relinquish, or repair it. Each layer of the paper ends on the governance question its own literature leaves open.
